\chapter{Research Methodology}
\label{ch:methodology}

This chapter details the experimental design employed to investigate cold-start latency reduction in AWS Lambda. The study adopts a controlled empirical approach with pre-registered hypotheses, explicit variable isolation, and rigorous statistical analysis following \citet{wen2025variance}'s recommendations for serverless performance characterization.

\section{Research Design}

This investigation employs a \textbf{controlled experiment} with a multi-factorial design. The independent variables (runtime, package size, memory allocation, warming strategy) are systematically manipulated while holding the workload constant. The primary dependent variable is \textit{Init Duration} (milliseconds), extracted from AWS CloudWatch REPORT log lines. Secondary metrics include total Duration, Billed Duration, and cost per 1,000 invocations.

\subsection{Experimental Phases}

The experiment proceeds through six phases, illustrated in Figure~\ref{fig:experiment_phases}:

\begin{enumerate}
    \item \textbf{Phase 0 (Packaging):} Build six deployment variants (3 runtimes × 2 package sizes) with identical SHA-256 computation workload. Verify digest uniformity via CI/CD pipeline.
    
    \item \textbf{Phase P (Pilot):} Execute small-sample runs (N=30 per configuration) under idle conditions (10-minute inter-invocation gaps) to estimate effect sizes and variance. Apply power analysis to determine required sample sizes for 80\% power at $\alpha = 0.05$.
    
    \item \textbf{Phase A (Baseline):} Collect multi-arm baseline data across all 6 function variants (without warmer) under idle, steady (one invocation per minute), and burst (100 concurrent invocations) traffic patterns.
    
    \item \textbf{Phase H1–H4 (Hypothesis Tests):} Execute targeted invocation campaigns for each hypothesis:
    \begin{itemize}
        \item H1 (Runtime): Compare minimal-package configurations across runtimes
        \item H2 (Package Size): Compare bloated vs. minimal within each runtime
        \item H3 (Warming): Enable EventBridge warmer (5-minute schedule), measure cold-start frequency
        \item H4 (Memory, Exploratory): Sweep memory allocations (128, 512, 1024, 1769, 3008 MB)
    \end{itemize}
    
    \item \textbf{Phase C (Collection):} Retrieve CloudWatch Logs via Logs Insights API with automatic pagination and retry. Parse REPORT lines to extract Init Duration, Duration, Billed Duration, Memory Used.
    
    \item \textbf{Phase S (Statistical Analysis):} Apply pre-registered analysis plan: normality tests (Shapiro-Wilk), hypothesis tests (parametric or non-parametric), multiple comparison correction (Holm-Bonferroni), effect size computation.
\end{enumerate}

% Placeholder for Figure: Experiment Phase Diagram
% Figure would show: Phase 0 → Phase P → Phase A → Phases H1-H4 → Phase C → Phase S

\subsection{Controlled Variables}

\textbf{Independent Variables (Manipulated):}

\begin{itemize}
    \item \textbf{Runtime:} Python 3.12, Node.js 20, Java 21 (AWS Lambda production runtimes as of September 2026)
    \item \textbf{Package Size:} Bloated (includes popular but unused dependencies) vs. Minimal (standard library or zero dependencies)
    \item \textbf{Memory Allocation:} 128 MB to 3008 MB in logarithmic steps
    \item \textbf{Warming Strategy:} None (baseline) vs. EventBridge rule (5-minute fixed schedule)
\end{itemize}

\textbf{Held Constant:}

\begin{itemize}
    \item \textbf{Workload:} Identical SHA-256 computation (20,000 rounds over fixed payload) across all variants, verified by digest equality
    \item \textbf{Architecture:} ARM64 Graviton2 (cost-effective per \citet{bluemke2025ijet})
    \item \textbf{Region:} Single AWS region (to eliminate geographic variance)
    \item \textbf{Invocation Mode:} Synchronous RequestResponse (AWS SDK invoke with wait)
    \item \textbf{Payload:} Fixed 1 KB JSON (\texttt{payloads/fixed\_payload.json})
    \item \textbf{Account:} Student's personal AWS account (ethics Scenario 3)
\end{itemize}

\textbf{Dependent Variables (Measured):}

\begin{itemize}
    \item \textbf{Init Duration} (ms): Primary outcome, extracted from CloudWatch REPORT lines when present (indicates cold start)
    \item \textbf{Duration} (ms): Total execution time including initialization and handler
    \item \textbf{Billed Duration} (ms): Rounded up to nearest 100 ms, used for cost calculation
    \item \textbf{Cold Start Indicator:} Binary (1 if Init Duration present, 0 otherwise)
    \item \textbf{Cost per 1,000 Invocations:} Computed via pricing formula (Section~\ref{sec:cost_model})
\end{itemize}

\section{Measurement Instrumentation}

\subsection{Init Duration Extraction}

AWS Lambda logs execution metrics to CloudWatch in REPORT lines with the following format when a cold start occurs:

\begin{verbatim}
REPORT RequestId: abc123 Duration: 450.23 ms 
Billed Duration: 500 ms Memory Size: 512 MB 
Max Memory Used: 128 MB Init Duration: 380.12 ms
\end{verbatim}

When the container is warm (reused from a previous invocation), the Init Duration field is absent:

\begin{verbatim}
REPORT RequestId: def456 Duration: 8.45 ms 
Billed Duration: 100 ms Memory Size: 512 MB 
Max Memory Used: 128 MB
\end{verbatim}

The log parser (\texttt{scripts/parse\_report\_metrics.py}) uses regular expressions to extract these fields, tagging each invocation as cold (Init Duration present) or warm (absent). This follows AWS's official cold-start definition: \textit{an invocation is cold if and only if its REPORT line includes Init Duration} \citep{bluemke2025ijet}.

\subsection{Log Collection Pipeline}

The collection pipeline (Figure~\ref{fig:pipeline_diagram}) operates as follows:

\begin{enumerate}
    \item \textbf{Invocation:} Python scripts (\texttt{invoke\_idle.py}, \texttt{invoke\_steady.py}, \texttt{invoke\_burst.py}) call Lambda functions via AWS SDK (boto3) with \texttt{InvocationType='RequestResponse'}. Each invocation records: request ID, client-side timestamp, function name, and cold-start intent (idle invocations expect cold, steady invocations expect warm).
    
    \item \textbf{Wait Period:} After invocation campaign completes, wait 5 minutes for CloudWatch log ingestion (AWS CloudWatch Logs Insights has 1–3 minute latency).
    
    \item \textbf{Query:} \texttt{collect\_logs.py} issues CloudWatch Logs Insights queries per function log group:
    \begin{verbatim}
    fields @timestamp, @message | 
    filter @message like /REPORT RequestId/
    \end{verbatim}
    Automatic pagination retrieves all REPORT lines.
    
    \item \textbf{Parsing:} \texttt{parse\_report\_metrics.py} extracts metrics into structured CSV with columns: \texttt{request\_id}, \texttt{function\_name}, \texttt{timestamp}, \texttt{duration}, \texttt{billed\_duration}, \texttt{memory\_size}, \texttt{max\_memory\_used}, \texttt{init\_duration} (nullable), \texttt{is\_cold} (boolean), \texttt{data\_mode} (\texttt{live}, \texttt{mock}, or \texttt{proxy}).
    
    \item \textbf{Validation:} Intended-cold invocations (10+ minutes since last call) that return warm are flagged and discarded from cold-start analysis. This prevents container recycling anomalies from biasing results.
\end{enumerate}

% Placeholder for Figure: Data Collection Pipeline Diagram
% Figure would show: Invoke → Wait → Query CloudWatch → Parse REPORT → Validate → CSV

\subsection{Cost Modeling}
\label{sec:cost_model}

AWS Lambda pricing (as of September 2026, US East region) consists of two components:

\begin{enumerate}
    \item \textbf{Request Charge:} \$0.20 per 1 million requests
    \item \textbf{Duration Charge:} \$0.0000166667 per GB-second (memory allocated in GB × billed duration in seconds)
\end{enumerate}

The cost per invocation is computed as:

\begin{equation}
C_{\text{invocation}} = 0.0000002 + \left(\frac{M}{1024} \times \frac{D_{\text{billed}}}{1000}\right) \times 0.0000166667
\end{equation}

where $M$ is memory in MB and $D_{\text{billed}}$ is billed duration in milliseconds.

Cost per 1,000 invocations is:

\begin{equation}
C_{1000} = 1000 \times C_{\text{invocation}}
\end{equation}

The decision matrix expresses each optimization as:

\begin{equation}
\text{ROI} = \frac{\Delta t_{\text{median}}}{\Delta C_{1000}}
\end{equation}

where $\Delta t_{\text{median}}$ is median latency reduction (ms) and $\Delta C_{1000}$ is cost increase per 1,000 invocations (USD).

\section{Statistical Analysis Plan}

All hypothesis tests are pre-registered in \texttt{docs/ANALYSIS\_PLAN.md} following \citet{wen2025variance}'s variance-aware methodology.

\subsection{Test Selection Criteria}

Figure~\ref{fig:stats_decision_tree} illustrates the statistical decision tree:

\begin{enumerate}
    \item \textbf{Normality Check:} Apply Shapiro-Wilk test ($\alpha = 0.05$) to each sample (Init Duration distributions per configuration).
    
    \item \textbf{Parametric Tests (if normal):}
    \begin{itemize}
        \item H1 (multi-group): One-way ANOVA with Tukey HSD post-hoc
        \item H2 (two-group): Welch t-test (unequal variance assumption)
        \item H3 (paired): Paired t-test on 20-minute block means
        \item H4 (multi-group): One-way ANOVA with Tukey HSD post-hoc
    \end{itemize}
    
    \item \textbf{Non-Parametric Tests (if non-normal):}
    \begin{itemize}
        \item H1 (multi-group): Kruskal-Wallis H test with Dunn post-hoc
        \item H2 (two-group): Mann-Whitney U test
        \item H3 (paired): Wilcoxon signed-rank test on block means
        \item H4 (multi-group): Kruskal-Wallis H test with Dunn post-hoc
    \end{itemize}
\end{enumerate}

% Placeholder for Figure: Statistical Decision Tree
% Figure would show: Shapiro-Wilk → Normal? → Yes: ANOVA/t-test, No: Kruskal-Wallis/Mann-Whitney

\subsection{Multiple Comparison Correction}

Holm-Bonferroni sequential correction is applied over five primary hypothesis tests:

\begin{enumerate}
    \item H1 (Runtime comparison)
    \item H2-Python (Package size, Python)
    \item H2-Node.js (Package size, Node.js)
    \item H2-Java (Package size, Java)
    \item H3 (Warming effect)
\end{enumerate}

H4 (Memory, exploratory) is not included in the family; its p-values are reported uncorrected with explicit "exploratory" designation.

\subsection{Effect Size Reporting}

All hypothesis tests report standardized effect sizes:

\begin{itemize}
    \item \textbf{Kruskal-Wallis:} Epsilon-squared ($\epsilon^2$, analogous to eta-squared)
    \item \textbf{Mann-Whitney U:} Rank-biserial correlation ($r_{\text{rb}}$, range $[-1, 1]$)
    \item \textbf{Welch t-test:} Cohen's d with Hedges correction for small samples
    \item \textbf{One-way ANOVA:} Eta-squared ($\eta^2$) or partial eta-squared
\end{itemize}

Effect sizes are interpreted using Cohen's conventional thresholds: small (0.01–0.06), medium (0.06–0.14), large ($> 0.14$) for epsilon-squared/eta-squared; small (0.1–0.3), medium (0.3–0.5), large ($> 0.5$) for rank-biserial.

\subsection{Power Analysis}

Prior to the main campaign, a pilot study (Phase P) collects $N = 30$ samples per configuration. Power analysis (\texttt{scripts/pilot\_power.py}) computes required sample sizes for 80\% power at $\alpha = 0.05$ given observed effect sizes and variance. Sample sizes are frozen before Phase H1–H4 to prevent data-dependent stopping (p-hacking).

\section{Ethics and Responsible Conduct}

\subsection{Ethics Scenario Classification}

This research falls under \textbf{Scenario 3} (self-owned infrastructure) of the NCI ethics framework:

\begin{itemize}
    \item No human participants (synthetic workload only)
    \item No secondary personal datasets
    \item Target system is the student's personal AWS account
\end{itemize}

A Declaration of Ethics Consideration is filed with the supervisor and retained for Moodle submission.

\subsection{Cost and Energy Awareness}

\textbf{Budget Guard:} A pre-invocation cost estimator (\texttt{scripts/budget\_guard.py}) halts execution if projected spend exceeds a configured daily cap (\$5 by default). Pilot-driven sample-size determination ensures the campaign remains within free-tier eligibility or low-cost thresholds.

\textbf{Energy Reporting:} While direct energy consumption is not measurable (AWS abstracts physical infrastructure), the study reports total GB-seconds and compute-time, which correlate with energy usage. ARM64 Graviton2 is selected for energy efficiency (20\% lower power consumption than x86\_64 per AWS documentation).

\subsection{Reproducibility and Data Provenance}

\textbf{Data Mode Segregation:} Every data point carries a \texttt{data\_mode} tag (\texttt{live}, \texttt{mock}, \texttt{proxy}). The analysis pipeline (\texttt{scripts/analyse.py}) refuses to mix modes, preventing accidental conflation of real and synthetic results.

\textbf{Proxy Benchmarks:} Local process-start measurements (Section~\ref{sec:proxy}) are always labeled "proxy" or "local init benchmark," never claimed as Lambda Init Duration. They serve to validate hypothesis directions before AWS execution.

\textbf{Mock Pipeline:} Synthetic data generated via \texttt{configs/mock\_model.yaml} validates end-to-end mechanics (invoke → collect → parse → analyze → figures). All mock outputs are watermarked "SYNTHETIC" and excluded from conclusions.

\section{Threats to Validity}

\subsection{Internal Validity}

\begin{itemize}
    \item \textbf{Container Reuse Anomalies:} AWS may reuse containers non-deterministically. Mitigation: 10-minute idle periods between intended-cold invocations; warm-returning invocations are flagged and discarded.
    \item \textbf{Platform Variance:} Host machine load, network latency to log collectors, micro-VM placement vary. Mitigation: Large sample sizes (determined by pilot), repeated measurements, non-parametric tests robust to outliers.
\end{itemize}

\subsection{External Validity}

\begin{itemize}
    \item \textbf{Workload Generalizability:} Results are specific to CPU-bound SHA-256 computation. I/O-bound workloads (S3 access, database queries) may exhibit different init-cost profiles.
    \item \textbf{Temporal Generalizability:} AWS Lambda platform evolves. Results reflect September 2026 platform state; future runtime optimizations may change magnitudes.
\end{itemize}

\subsection{Construct Validity}

\begin{itemize}
    \item \textbf{Init Duration Definition:} AWS does not publicly document Init Duration composition (micro-VM vs. runtime vs. application). This study treats it as an atomic metric.
\end{itemize}

\section{Summary}

This chapter detailed the controlled experimental design, measurement instrumentation, statistical analysis plan, and ethics considerations. The methodology emphasizes reproducibility (data provenance tracking, pre-registered hypotheses), rigor (power analysis, effect sizes, multiple comparison correction), and honesty (explicit labeling of proxy and mock data).

The next chapter presents the design specifications of the Lambda functions, Infrastructure-as-Code template, and workload implementation.